% l1_content.tex
% Частина конспекту лекції "Python. Вступ та синтаксис", оформлена відповідно до вимог ОНУ.





\chapter{\parbox{0.85\textwidth}{Чому Python?  Створення та запуск файлів. Робота в інтерпретаторі.}}

\section{Історія та Глобальний вплив}

Створення Python розпочав Гвідо ван Россум (Guido van Rossum) у 1991 році під час роботи над розподіленою ОС «Амеба». Йому була потрібна розширювана мова, яка б забезпечувала підтримку системних викликів. За основу було взято мови ABC та Modula-3. Назву Гвідо обрав на честь популярного комедійного шоу BBC «Літаючий цирк Монті-Пайтона», а не на честь змії. Відтоді Python активно розвивався за підтримки організацій, у яких працював його автор. Сьогодні мова вдосконалюється ще швидше, адже над нею працює не лише команда творців, а й велика міжнародна спільнота програмістів. Втім, остаточне слово щодо напряму розвитку Python і надалі залишається за Гвідо ван Россумом.

\subsection*{Історія розвитку Python}

\begin{itemize}
	\item \textbf{Python 2.0} --- випущений у 2000 році. Серед нововведень: спискові включення, підтримка Unicode, вдосконалений збирач сміття.
	\item \textbf{Python 3.0} --- випущений у 2008 році. Основна особливість --- несумісність із попередніми версіями, що вимагало адаптації бібліотек.
	\item \textbf{Python 2.7} --- останній реліз гілки 2.x. Його підтримку офіційно завершено у 2020 році, що ознаменувало перехід до епохи Python~3.
\end{itemize}

\subsection*{Глобальний вплив та переваги}

Python став однією з найпопулярніших мов програмування у світі завдяки простоті синтаксису та універсальності. Його активно використовують такі компанії, як \textbf{Google, Netflix, Facebook, Instagram, Spotify} та багато інших.

\medskip
Основні сфери застосування:
\begin{itemize}
	\item \textbf{Наука та дослідження} --- бібліотеки NumPy, SciPy, pandas.
	\item \textbf{Машинне навчання та штучний інтелект} --- TensorFlow, PyTorch, scikit-learn.
	\item \textbf{Веб-розробка} --- фреймворки Django, Flask.
	\item \textbf{Автоматизація та DevOps} --- скрипти для адміністрування, тестування, CI/CD.
	\item \textbf{Аналіз даних та візуалізація} --- matplotlib, seaborn, Jupyter Notebooks.
\end{itemize}

\medskip
Таким чином, Python використовується практично всюди --- від навчальних проєктів до масштабних систем світових технологічних гігантів.



\section{Переваги мови або чому Python}

\subsection*{Основні переваги}

\begin{itemize}
	\item \textbf{Якість програмного забезпечення.} Зручність читання, ясність і висока якість відрізняють Python від інших мов. Єдиноманітність оформлення коду полегшує його розуміння навіть для тих, хто не брав участі у створенні програми.
	
	\item \textbf{Висока швидкість розробки.} У порівнянні з C, C++ та Java, Python значно підвищує продуктивність праці розробника. Обсяг коду зазвичай становить третину або навіть п’яту частину еквівалентного коду на C++ чи Java. Програми запускаються відразу, минаючи етапи компіляції та зв’язування.
	
	\item \textbf{Перенесення програмного коду без витрат.} Більшість програм виконується без змін на основних платформах. Перенесення з Linux у Windows часто зводиться до копіювання файлів. Python також підтримує створення графічних інтерфейсів, доступ до баз даних, веб-додатків та інших типів програм.
	
	\item \textbf{Бібліотеки підтримки.} У складі Python є велика стандартна бібліотека, що охоплює пошук тексту, роботу з мережею та інші функції. Крім того, існують численні сторонні бібліотеки: для веб-розробки, математичних обчислень, ігор тощо. Наприклад, \texttt{NumPy} --- потужний інструмент для наукових обчислень.
	
	\item \textbf{Інтеграція компонентів.} Python легко взаємодіє з іншими мовами та платформами: може викликати функції з бібліотек на C/C++, інтегруватися з Java, взаємодіяти з COM та .NET, працювати через протоколи SOAP, XML-RPC, CORBA.
	
	\item \textbf{Вільне використання.} Python можна використовувати та розповсюджувати абсолютно безкоштовно.
\end{itemize}

\subsection*{Класифікація мови}

\begin{enumerate}
	\item \textbf{Високого рівня.} Код легко читати і писати, він схожий на англійську мову.
	\item \textbf{Загального призначення.} Python вирішує широкий спектр завдань.
	\item \textbf{Інтерпретована.} Для виконання коду не потрібна компіляція.
\end{enumerate}

Python часто позиціонують як \textbf{об’єктно-орієнтовану мову сценаріїв}, що дозволяє створювати прості програми у вигляді одного файлу, використовуючи об’єкти.

\subsection*{Недоліки}

\begin{itemize}
	\item \textbf{Швидкодія.} Програми на Python можуть працювати повільніше за аналоги, написані компілюючими мовами (наприклад, C).
	\item \textbf{Байт-код.} Python перетворює вихідний код у байт-код, який інтерпретується. Це забезпечує переносимість програм, але відсутність машинного коду може знижувати продуктивність.
\end{itemize}

\section{Середовище для роботи та Інтерпретатор}


Python --- це не лише мова програмування, а й програмний пакет, що включає \textbf{інтерпретатор} --- модуль, який створює та виконує інші програми.

\subsection*{Онлайн середовища розробки}

Сьогодні є можливість писати код на Python у веб-сервісах, які дозволяють запускати програми прямо у браузері, без встановлення додаткового програмного забезпечення:
\begin{itemize}
	\item \url{https://replit.com/~}
	\item \url{https://ideone.com/}
	\item \url{https://www.online-python.com/}
\end{itemize}

\medskip
У цих сервісах достатньо створити акаунт, написати код та запустити його онлайн.

\subsection*{Командний рядок інтерпретатора}

Найпростіший спосіб писати й запускати програми на Python --- вводити інструкції безпосередньо у \textbf{командний рядок інтерпретатора}. Це дозволяє швидко перевіряти окремі команди та експериментувати з кодом.

\subsection*{Локальне встановлення Python}

Щоб працювати локально, потрібно завантажити Python з офіційного сайту:
\begin{itemize}
	\item \url{https://www.python.org/downloads/}
\end{itemize}

Після встановлення пакету Python створюється інтерпретатор та бібліотека підтримки. Також доступне інтегроване середовище розробки \textbf{IDLE}, яке встановлюється разом із Python і дозволяє писати, запускати та відлагоджувати програми у зручному графічному інтерфейсі.

\section{Виконання вказівок у Python}

\subsection*{Інтерактивний режим}

В інтерактивному режимі користувач вводить інструкції та натискає \texttt{Enter}, після чого результат відображається рядком нижче.

\begin{verbatim}
	>>> 2 + 2
	4
	>>> "Hello"
	'Hello'
\end{verbatim}

Інтерактивний режим дозволяє:
\begin{itemize}
	\item швидко перевіряти фрагменти коду;
	\item проводити експерименти без створення файлів;
	\item одразу виявляти помилки.
\end{itemize}

Приклад помилки:
\begin{verbatim}
	>>> print(y)
	Traceback (most recent call last):
	File "<pyshell#15>", line 1, in <module>
	print(y)
	NameError: name 'y' is not defined
\end{verbatim}

\subsection*{Виведення повідомлень на екран}

Для виведення інформації використовується команда \texttt{print}:

\begin{verbatim}
	>>> print("Hello World")
	Hello World
	>>> print(2 + 2)
	4
\end{verbatim}

У інтерактивному режимі можна вводити вирази без \texttt{print}, але у файлах інструкцію \texttt{print} потрібно писати обов’язково.

\subsection*{Програми у файлах}

Програму на Python можна написати у будь-якому текстовому редакторі та зберегти у файл з розширенням \texttt{.py}.  
Наприклад, створимо файл \texttt{hello.py}:

\begin{verbatim}
	print("Hello world!")
	print("Привіт світ")
	input()
\end{verbatim}

Запуск програми здійснюється подвійним клацанням по файлу.  
Такі програми виконуються у консольному вікні («чорному вікні»), де введення здійснюється з клавіатури. Якщо у програмі є помилки, у консолі буде показано повідомлення про них.

\medskip
Таким чином, Python дозволяє працювати як у інтерактивному режимі, так і створювати повноцінні програми у файлах.

\section{Перевірка коду на помилки в IDLE}

\subsection*{Запуск програми у IDLE}

\begin{enumerate}
	\item Запускаємо середовище \textbf{IDLE}.
	\item У меню \textbf{File} обираємо пункт \textbf{Open} та у вікні шукаємо файл \texttt{hello.py}.
	\item Відкриється нове вікно з текстом програми. У цьому вікні у меню \textbf{Run} вибираємо \textbf{Run Module} або натискаємо клавішу \texttt{F5}.
\end{enumerate}

\subsection*{Результати виконання}

\begin{itemize}
	\item Якщо у програмі немає помилок, введення та виведення здійснюється у вікні інтерпретатора.
	\item Повідомлення про помилки також відображаються у \textbf{IDLE}.
	\item Якщо програма працює некоректно, але не видає помилок, для пошуку проблем можна вставити додаткові інструкції \texttt{print}.
\end{itemize}

\subsection*{Приклад використання \texttt{print} для відлагодження}

\begin{verbatim}
	print("Значення змінної x:", x)
\end{verbatim}

Це допоможе зрозуміти, які значення мають змінні під час виконання програми.

\section*{Контрольні питання}

\begin{enumerate}
	\item Які основні сфери застосування Python ви знаєте?
	\item Назвіть головні переваги мови Python.
	\item Які недоліки Python були згадані у тексті?
	\item Як класифікується мова Python (за рівнем, призначенням, способом виконання)?
	\item Що таке інтерпретатор Python і яку роль він виконує?
	\item Які онлайн‑середовища можна використати для написання та запуску програм на Python?
	\item Як працювати з Python у командному рядку інтерпретатора?
	\item Які кроки потрібно виконати для локального встановлення Python?
	\item Чим відрізняється інтерактивний режим від запуску програм у файлах?
	\item Як у Python здійснюється виведення повідомлень на екран?
	\item Як перевіряти код на помилки у середовищі IDLE?
	\item Які результати виконання програми можна побачити у IDLE?

\end{enumerate}
